\chapter{Introduction}

\vspace{-1cm}
\section{Initial motivation}
\label{sec:initialMotivation}

Unexplored simplicity. This was the reason for the birth of this work.
Let me explain. Topology deals with, among other objects, the so
called {\em 3-manifolds}. A basic type of 3-manifold is a {\em
closed, connected, oriented 3-manifold} and, in this work, the term
{\em space} will be used as a synonym for it.\footnote{The term
``space'' is also used as a synonym for a 3-manifold (any
3-manifold), but here we will use it as a synonym only for a closed,
connected, oriented 3-manifold.} In 1994, Kauffman and Lins
\cite{KauffmanAndLins1994} introduced a new way to present spaces.
They named their new presentation {\em blinks}. In fact, on their
work, blinks were appreciated as being a new concise presentation
for spaces, but its main importance there was its use as an
intermediate step to convert a {\em blackboard framed link}
presentation of a space into a {\em 3-gem} presentation of the same
space. This conversion, among other results, is central to this work,
but we will get to it later. What I want to say now is what
attracted me to this object named blink. Here it is.
\vspace{-0.3cm}
{\begin{center} \it
Blink states that a triangle is a space. That a square is a space.
That any plane graph (\ie any drawing in a plane made of
points connected by curves that do not intersect) is a space.
\end{center}}
\vspace{-0.3cm}
\noindent What space has the form of a
triangle? And the form of a square? I found this connection
between plane graphs and spaces very elegant. This elegance
becomes even more special when we see that other space
presentations have more ``complicated'' drawings than
blink presentation does. For example, the following three
drawings are different presentations for space $\IS^3\!/\!<3,3,2>$. \footnote{That is, the quotient of the $3$-sphere
under the action of the binary tetrahedral subgroup, which is a non-abelian group of order 12.}
\begin{center}
\includegraphics[width=8cm]{fig/presentations.eps}
\end{center}
The first is a blink presentation, the second is a blackboard
framed link presentation and the third is a 3-gem presentation.
In this example it is clear that the blink presentation is
simpler. Its perceptual complexity is smaller. Indeed this is
always the case. No other space presentation has simpler drawings
than blink does. Blackboard framed links don't, 3-gems don't.
{\it special spines} don't, {\it Heegaard diagrams} don't.

This simplicity aspect of blinks allied to the fact that they were
not studied before, except as a by product on the books
\cite{KauffmanAndLins1994, Lins1995} made me decide to explore it.
To see the greatest number of spaces through blink
drawings was the initial intent, but the hope was also that
so much elegance and simplicity is not in vain and it could,
by the fact that it was not explored before, hold
some yet unknown secret of spaces.

\section{Historical overview}
\label{sec:historicalOverview}

In 1962 Lickorish \cite{Lickorish1962} revolutionized the area of
spaces by proving in a purely combinatorial way a result first
published by Wallace \cite{Wallace1960} two years before by means of
differential topology. The result has, as a corollary, the following
fact:

{\em Given any space $M$, there exists a finite number $k$ of
disjoint solid tori $T_i\subset M$ such that $M \backslash
\bigcup_{i=1}^k T_i$ is homeomorphic to
$\IS^3\backslash\bigcup_{i=1}^k T'_i$, for a corresponding set of
disjoint solid tori $T'_i \subset \IS^3$.}

Another consequence of Lickorish's Theorem permits the
presentation of an arbitrary space as a link diagram in the plane
with an integer attached to each one of the components, the so
called {\em framed links}. The revolution was completed in 1978 when
Kirby \cite{Kirby1978} published his now famous calculus on framed links.
This paper spurred an enormous activity in the
area paving the way to prove deep theorems by a specific kind of
diagrammatic calculus with well defined rules.

Framed links can be freed from the integers associated to their
components if one introduces curls in their projections so as to
have the frame of a component equal to its writhe, thus producing a
{\em blackboard framed link} or a {\em BFL}. In this way, every
space presents itself as a blackboard framed link. A reformulation
of Kirby's calculus into blackboard framed link language is presented
in \cite{Kauffman1991}: a formal BFL calculus. The importance of the
BFL presentation is
testified by the fact that we can obtain from
it the Witten-Reshetikhin-Turaev quantum invariant (or WRT-invariant) of
the induced space \cite{Witten1989,ReshetikhinAndTuraev1991}. This is
not possible, at least it is unknown at the present, from a
triangulation \cite{Hempel1976}, from a Heegaard diagram
\cite{Rolfsen1976} or from a {\em special spine} \cite{Matveev2003}
of the space. In this work the {\em QI} or quantum invariant of a
space mean their WRT-invariants computed at $A=e^{i\pi/2r}$
\cite{KauffmanAndLins1994}.

A BFL can be reformulated into the object named {\it blink} also
introduced in \cite{KauffmanAndLins1994}. The presentation by blink
is concise and permits, as BFLs does, the computation of best invariants
available. But how do we prove that two blinks yielding the same
invariants are manifestations of the same space? Kirby's moves, although a
theoretical masterpiece, are in this practical context of no use,
except in very limited circumstances. For this task, the TS-moves
and U-move on gems of Lins \cite{Lins1995} are much better. The
applicability of the gem simplifying dynamics is available because
from a blink it is straightforward to produce a gem inducing the
same space \cite{KauffmanAndLins1994}. Summarizing, the task of
classifying blinks has to rely on gem theory which is, at present,
an indispensable complement to perform the task. The topological
classification of gems with 30 vertices (extending the classification
of \cite{Lins1995} using again TS- and U-moves) has been recently
accomplished~\cite{CasaliCrist2006}.

\newpage

\section{What we did}

In Section~\ref{sec:initialMotivation} we said our intent was to see
the greatest number of distinct spaces through blinks (\ie blink drawings).
To be able to accomplish this task, the most important question we must know
how to answer is whether two blinks $A$ and $B$ are
actually presentations of the same space. A fact we did
not mention before is that any space has infinitely
many distinct blink presentations. So, to answer this question is
not just checking that $A = B$ or $A \neq B$. The relation
space-blink is not a one to one relation.

As mentioned in Section~\ref{sec:historicalOverview}, blink is a direct
reformulation of blackboard framed link (BFL). It is easy to go from a blink
presentation of a space to a BFL presentation of the same space and
vice-versa. We also mentioned that a formal calculus for BFLs is well
known. This
means that a set of BFL operations or moves is known such that any pair
of BFLs induce the same space if and only if there
is a path, \ie a finite sequence of operations or moves in this set, that
transforms one BFL into the other. So, we could answer our question like
this: obtain a BFL presentation for $A$, obtain a BFL presentation for
$B$, then show a path in the BFL calculus transforming $A$ into $B$, proving
that they induce the same space; or show that such a path does not exist, and
conclude that they induce different spaces. Although this approach is correct and theoretically
possible, it is not practical. How to show a proof that a path does not
exist? Despite of this practical gap, one of the contributions of
this thesis is a reformulation of the BFL calculus into a purely blink
calculus.% But this also did not help in answering our practical question.

Section~\ref{sec:historicalOverview} also mentions that from a blink we
are able to calculate some space invariants. For example, it is possible
to obtain the homology group or the Witten-Reshetikhin-Turaev
quantum invariant of its space. Indeed, this is the first thing we do
to answer whether blink $A$ induces the same space as blink $B$. We
calculate these two invariants and if any of them are different
we can answer for sure that the blinks induce different spaces. But,
if they are both the same, we cannot say their spaces are the same.
No complete invariant for spaces is known. So, any known space
invariant may fail to distinguish different spaces.

When the space of blink $A$ is not distinguished from the space
of blink $B$ by the space invariants, we must use another tool to answer
our question. This other tool is 3-gem theory. The book
\cite{KauffmanAndLins1994} shows a way to obtain a 3-gem presentation
from a blink presentation inducing the same space. We improved this
algorithm in Chapter~\ref{chap:gems}. In \cite{Lins1995}
a nice algorithm to simplify 3-gems is presented. So the last thing
we do to answer our question is to check whether the
blinks $A$ and $B$ not distinguished by the space invariants have
their 3-gems simplified to a common 3-gem. If this is the case then
we are sure that $A$ and $B$ induce the same space. If not, then we
are not sure. It is a hint that they are different but this cannot
be said. For small blinks, as we will see later, there are only two
uncertainties left out of $\approx 500$.

This approach of testing the homology group and the WRT-quantum invariant
to distinguish blinks and then, if not distinguished,
applying the 3-gem simplifying algorithm to show that they induce
the same space was very successful in our experiments as we will
see in Chapter~\ref{chap:census}. Its only constraint is that
it works only for small blinks and 3-gems. The computational
effort to calculate quantum invariants or to simplify 3-gems is exponential
in the sizes of the blinks and of the 3-gems, respectively.

Let's return to our initial intent: isolating the largest number of
spaces through blinks. We already know how to test if two (small)
blinks do induce the same space or not. The next important thing
to define is for what blinks we are going to ask these questions.
To try all possible blinks is prohibitive and unnecessary. As we
will see, we can search for spaces in only a small fraction of all possible
blinks and yet not lose anything. To get to this optimization
we first developed a useful decomposition/composition theory
of blinks in Chapter~\ref{chap:blinks} (actually the theory was developed
for its combinatorial counterpart: the g-blink). Then, using the
results and operations of blink calculus and BFL calculus we filtered
some redundant blinks. This resulted in a small set of blinks for which we
could identify all spaces. We also could isolate interesting sets
of small blinks inducing the same spaces such that a path in blink
calculus or BFL calculus is not trivial to identify. These sets may
lead to new ideas for theorems or space invariants.

A contribution of this thesis is a computer program named \textsc{Blink}.
It was responsible for most of the figures in this document. It also
supports the most important concepts discussed in the following
chapters and we hope that it will become popular and help researchers and
students to learn, do research or just appreciate spaces through
the language of blinks, blackboard framed links or 3-gems which are
the objects that it supports at the moment.

\section{The structure of this thesis}

In Chapter~\ref{chap:topologyEtc} we begin with a
review of the basic topological language and concepts needed in the remaining of
the thesis. We then introduce {\it knots and links} and their diagrams.
After that we introduce {\it framed links} and
{\it blackboard framed links} (BFL) and show how they encode spaces.
A calculus for blackboard framed links is presented.

In Chapter~\ref{chap:blinks} we define the motivating object of this
work: a {\it blink}. We show that blinks are a simple
reformulation of blackboard framed links with the advantage of having
simple drawings. We then reformulate the BFL calculus shown in
Chapter~\ref{chap:topologyEtc} in blink language. From
blinks we define a new combinatorial object named {\it g-blinks}. We show
how to obtain the homology group and quantum invariant of
Witten-Reshetikhin-Turaev invariants from a g-blink. The code
of a g-blink is then presented. Then some involutions
of g-blinks are defined: dual, reflection and refDual. The concept
of a representative g-blink is introduced based on the previous
results shown for g-blinks. We end the chapter showing how to
identify all spaces that are induced by small blinks and what
is the missing piece to do that: a way to prove homeomorphisms.

In Chapter~\ref{chap:gems} we define 3-gems: another way to present
spaces. We then show some moves that can be done in
3-gems without changing the induced space. These moves yields a
viable computational way to prove homeomorphism of spaces: a
combinatorial simplification dynamics of 3-gems. To connect blinks
and 3-gems we show an improved way to obtain, from a blink,
a 3-gem inducing the same space. We finish this chapter with the
proof, via 3-gems, of a theorem on g-blinks stated in
Chapter~\ref{chap:blinks}: the partial reflection theorem.

In Chapter~\ref{chap:census} we present the computational
experiments and results that we have obtained. We define formally
what we are searching: a census of prime spaces. Then we construct
a small set named $U$ that has the {\it 9-prime-unavoidable} property.
We then show how we topologically identified the spaces of every
g-blink in $U$. We finish the chapter exploring another set of
g-blinks: simple 3-connected monochromatic blinks.

In Chapter~\ref{chap:conclusion} we review the main contributions
of this work, talk a little about the program \textsc{Blink} and
about a theoretical contribution that we did not finish on
time to this thesis: a polynomial algorithm to obtain the blink
of a 3-gem. Some research problems that can be explored as a
continuation of this work are also shown.







